\section{The formalisation}\label{app:formal}

\paragraph{What was checked.}
The development is written in Lean 4 \citep{Lean4}, toolchain \lean{v4.34.0-rc2}, against the
Mathlib library \citep{Mathlib} at a fixed revision recorded in the repository manifest. It
comprises 115 source files, about 35,000 lines, and about 1,500 theorem declarations. A Lean proof
is a term whose type is the statement; the kernel checks that the term is well-formed. What the
reader is asked to trust is the kernel, the axioms, and that the formal statement means what the
prose says. The first two are shared with every Lean development. For the third, each result in the
text names its declaration, and Section~\ref{sec:model} describes the formal objects in enough
detail that the statements can be read; the hypotheses in the text are the hypotheses in the code.

\paragraph{The household problem.}
The dynamic programme is built on a general stochastic dynamic-programming layer with
extended-real rewards, so that $u=-\infty$ at zero consumption is handled without a consumption
floor. The principle of optimality, existence and uniqueness of the bounded value function,
Berge's maximum theorem for the policy (which Mathlib lacks; ours is built on its hemicontinuity
definitions), and the Blackwell contraction are theorems of that layer. The asset region
$[\underline a,\bar a]$ is a type parameter, which is what lets the same theorems serve Aiyagari and
Huggett.

\paragraph{The rate as a state.}
Continuity of the policy in the interest rate is obtained by adjoining the rate to the state and
solving one dynamic programme on the product, so that joint continuity in the state is continuity
in the rate. This avoids any parametric-continuity argument, which would need a modulus that the
Bellman contraction does not obviously supply.

\paragraph{Secant slopes.}
Lipschitz bounds, monotonicity in assets, monotonicity in the rate and the Euler inequalities are
all proved by comparing plans, with no derivative of the value function. Where the development
proves the value function differentiable it uses the sandwich lemma of \citet{ClausenStrub2020},
formalised as \lean{DifferentiableSandwich}. The uniqueness proof does not depend on it.

\paragraph{Probability.}
Stationary distributions are Borel probability measures on a compact metric space, in Mathlib's
\lean{ProbabilityMeasure} type with the weak topology; the Krylov--Bogolyubov existence argument,
weak continuity of aggregation, and the Doeblin contraction on bounded continuous functions are
proved in \lean{Distribution/}. The normal distribution behind Tauchen's method is Mathlib's
\lean{gaussianReal}.

\paragraph{Reproducing the check.}
With Lean and Mathlib installed, \lean{lake build} at the repository root compiles the
development, and \lean{\#print axioms} on any declaration in the index reports the axioms it uses.
The repository is public and the history records the order in which the results were obtained.
